\section{Exact calculations}

All local formulas in Section~\ref{sec:local} are rational identities
in $r$, or polynomial identities after the substitution
$w=2-r-r^{-1}$. They require no Taylor truncation or numerical
approximation. The multiplication, pairing and cap data needed to
reproduce them are printed in \eqref{eq:M}--\eqref{eq:B}.
As one numerical specialization in exact arithmetic, at $r=2$
the characteristic coefficients in decreasing powers of $T$ are
\[
 1,\quad\frac{225}{32},\quad\frac{14027}{384},\quad
 \frac{158613}{4096},\qquad
 \tr(B^2)=-\frac{72557}{3072}.
\]
The general characteristic identity, rather than this single
specialization, is used in the proof.

The accompanying code gives two implementations: matrix reconstruction
in \path{checks/local_curve_degree_three.py} and quotient-algebra
inversion with a resultant in
\path{checks/degree_three_quotient_check.py}. Both use SymPy 1.14.0.
The first also checks degrees one and two and the antidiagonal
eigenvalues \eqref{eq:UV}. The script
\path{checks/integrality_counterexample.py} checks the connected
subtraction, the Laurent coefficient \eqref{eq:local-pt}, and all
56 unordered nontrivial effective decompositions of $(3,4)$.
Their minimum sum in Lemma~\ref{lem:split} is zero. The lemma's
three-case argument proves the inequality without this enumeration.

The toric data are printed in \eqref{eq:rays}--\eqref{eq:nef-matrix}.
The script \path{checks/integrality_counterexample_geometry.py}
checks the strict support function, anticanonical support, nef lifts
and lattice-basis test. It reads the same ray and cone lists from
\path{checks/x9_lift_bound.py}. The latter runs with SageMath and
checks completeness and smoothness of the fan. The contracted-face
and surface calculations are reproduced by
\path{checks/x9_contracted_support.py}, also using SageMath.
The integer degree table \eqref{eq:degrees} can be recovered from
the displayed divisors and surface fans.

These finite computations concern the local algebra and toric data.
The pure-curve filtration, moduli-scheme argument and virtual sign
are the geometric steps establishing that these data yield a
compact counterexample.
